\bookchapter{第十三章：工具类App推荐}

\chapterintro{本章精选效率、健康、出行类工具，并补充实用系统与跨设备能力：清浊等清理与防骚扰、手机自带拦截开关、天气与地震预警；向日葵/ToDesk 远控电脑、Trae 等 AI 办公、人工桌面、LocalSend 等局域网快传、跨设备剪贴板、Aegis/Microsoft Authenticator 动态口令、滴答清单日程与纸质扫描。}

\booksection{一、效率工具：管理你的时间与注意力}

效率工具是数字生活的“骨架”，帮助你把零散的想法、任务、资料整理得井井有条。下面按场景分类推荐。

\booksubsection{待办清单（任务管理）}

\begin{itemize}
  \item \textbf{滴答清单（TickTick）}适用人群：需要把工作、生活、学习任务与日程统一记录的人；本节重点推荐之一。核心卖点：跨手机/电脑/平板/Web 同步；\textbf{清单 + 日历 + 番茄钟}三合一；自然语言输入「明天下午3点开会」「每周一交周报」自动识别时间与重复规则；支持子任务、标签、优先级、附件与协作清单。适合「日程安排 + 任务打卡 + 每日回顾」一体化管理。
  \item \textbf{滴答清单怎么用得更顺}：
  \begin{itemize}
    \item \textbf{收集箱}：先全部丢进「今天/收集箱」，晚上再归类到工作、家庭、学习清单。
    \item \textbf{日历视图}：会议、截止日拖到日历上看冲突，比只看待办列表更直观。
    \item \textbf{重复任务}：吃药、缴费、健身设成规则任务，不用每次手填。
    \item \textbf{番茄钟}：在任务上直接开专注，统计本周「真正做完了多少」。
    \item \textbf{与手机功能联动}：语音助手说「提醒我…」可进系统提醒；重要截止日同时开系统日历双保险。
  \end{itemize}
  \item \textbf{Todoist}适用人群：追求极简、跨平台用户。核心卖点：界面克制清爽，项目层级清晰（任务子任务子子任务），karma 积分激励习惯养成。
  \item \textbf{Microsoft To Do}适用人群：办公党、Outlook用户。核心卖点：完全免费、无广告，"我的一天"功能每天清空、减轻心理负担，与微软生态深度集成。
  \item \textbf{系统便签/提醒事项}适用人群：不想装第三方App的轻度用户。核心卖点：原生体验、零延迟，Siri/小爱语音直接添加，与系统 widget 联动。
\end{itemize}

\booksubsection{番茄钟（专注力训练）}

\begin{itemize}
  \item \textbf{Forest}适用人群：管不住手、易刷手机的人。核心卖点：专注时种树，提前离开 App 树会枯死，用"虚拟生命"建立正向反馈，可与好友合种。
  \item \textbf{潮汐}适用人群：需要白噪音辅助专注/入睡的人。核心卖点：自然音效（雨声、海浪、森林）+番茄钟+冥想三合一，界面极简治愈。
  \item \textbf{原则}：经典 25 分钟专注 + 5 分钟休息，可根据自身状态调整为 50/10 或 90/20。
\end{itemize}

\booksubsection{日历（时间可视化）}

\begin{itemize}
  \item \textbf{系统日历}适用人群：所有用户。核心卖点：原生推送最稳、耗电极低，可与 Siri/小爱语音添加。
  \item \textbf{Google Calendar}适用人群：多设备用户、需要与他人共享日程者。核心卖点："目标"功能自动帮你安排健身/学习时间，与 Gmail 智能识别邮件中的事件。
\end{itemize}

\booksubsection{笔记类（知识沉淀）}

\begin{itemize}
  \item \textbf{Notion}适用人群：需要把笔记、任务、数据库、wiki 统一管理的人；学生党、自由职业者、团队协作。核心卖点：全能知识库，Block 模块化编辑，可搭建个人主页、阅读清单、项目看板，模板生态丰富。
  \item \textbf{Obsidian}适用人群：深度知识管理者、研究者、写作人。核心卖点：本地 Markdown 存储，双向链接 + 关系图谱，让你的笔记形成"第二大脑"，数据完全归你所有。
  \item \textbf{flomo}适用人群：想记录碎片灵感却不想写长文的人。核心卖点：卡片式笔记，类似发朋友圈一样记录想法，支持 \# 标签分类，按时间线回溯。
  \item \textbf{苹果备忘录}适用人群：iPhone/Mac 用户、轻度记录需求。核心卖点：启动快、支持扫描文档/手绘/清单/iCloud 同步，免费且与系统深度集成。
  \item \textbf{印象笔记}适用人群：需要大量网页剪藏、资料归档的资深用户。核心卖点：网页剪藏插件强大，OCR 识别图片文字，老牌稳定，但高级版需付费。
\end{itemize}

\booksubsection{PDF 工具}

\begin{itemize}
  \item \textbf{PDF Expert}适用人群：iPad/iPhone 用户、需要批注合同的商务人士、学生党。核心卖点：iOS 上公认最强的 PDF 编辑器，标注流畅、支持签名、合并拆分、填写表单。
  \item \textbf{WPS}适用人群：需要同时处理 Word/Excel/PPT/PDF 的办公族。核心卖点：一站式办公套件，国内网盘/微信文件兼容好，免费版功能已足够日常。
  \item \textbf{Adobe Acrobat}适用人群：对 PDF 格式规范要求高的人（法律、设计、出版）。核心卖点：官方标准、转换 PDF/A 格式无误，跨设备同步与电子签名功能专业。
\end{itemize}

\booksubsection{扫描类（纸质文件数字化）}

纸质合同、发票、作业、白板笔记、身份证复印件……手机摄像头已经能替代大部分扫描仪。核心流程：\textbf{拍清楚 自动裁边与矫正 导出 PDF/图片 OCR 可选}。

\begin{itemize}
  \item \textbf{扫描全能王（CamScanner）}适用人群：经常需要扫描合同、票据、证件的商务人士和学生。核心卖点：边缘自动识别、透视矫正、OCR 文字识别，可生成多页 PDF 并云同步。
  \item \textbf{微软 Lens}适用人群：学生党、需要把白板/书页扫描成 Word 的人。核心卖点：免费无广告，可一键导出到 OneNote/Word，识别文字准确率高。
  \item \textbf{苹果备忘录扫描}适用人群：iPhone 用户、轻度扫描需求。核心卖点：无需装 App，备忘录内点击相机图标即可扫描，自动去阴影、签名标注，免费好用。
  \item \textbf{白描 / 扫描全能王网页版 / 品牌相机「文档扫描」}：华为、小米、OPPO、vivo 系统相机大多内置「文档/扫描」模式；安卓也可用 Google 相机或系统扫一扫。适合应急拍一页合同。
  \item \textbf{拍照小技巧}：光线均匀、手机与纸面平行、对准四个角；多页文件按顺序拍完再合并 PDF；敏感证件注意水印与保存位置，用完可从「最近删除」彻底清理。
\end{itemize}

\booksection{二、健康运动}

健康类 App 是手机最被低估的“私人助理”——它们能持续追踪你的睡眠、心率、运动量，比偶尔去一次医院更能反映长期状态。

\booksubsection{运动健身}

\begin{itemize}
  \item \textbf{Keep}适用人群：居家训练者、零基础新手、想减脂塑形的人。核心卖点：海量免费课程，从 HIIT 到瑜伽全覆盖，AI 自动生成训练计划，社区氛围浓厚。
  \item \textbf{Nike Training Club（NTC）}适用人群：追求专业教练指导、希望系统训练的人。核心卖点：Nike 官方出品，课程全部免费，动作演示由专业教练录制，按肌群/器械分类清晰。
  \item \textbf{苹果健身（Fitness）}适用人群：Apple Watch 用户、想用闭环激励自己动起来的人。核心卖点：每日三个活动圆环（移动/锻炼/站立），可视化目标，与手表心率数据无缝打通。
\end{itemize}

\booksubsection{睡眠追踪}

\begin{itemize}
  \item \textbf{AutoSleep}适用人群：Apple Watch 用户、想知道自己每晚深睡/浅睡比例的人。核心卖点：戴表入睡自动追踪，无需手动启停，给出睡眠质量评分与就寝建议。
  \item \textbf{Sleep Cycle}适用人群：不戴手表睡觉、想用手机监测的人。核心卖点：通过麦克风分析翻身/呼噜声，在浅睡阶段智能唤醒，让你起床不痛苦。
  \item \textbf{小米运动健康}适用人群：小米手环/手表用户。核心卖点：免费且与小米手环深度联动，心率/血氧/睡眠一体追踪，性价比高。
\end{itemize}

\booksubsection{冥想放松}

\begin{itemize}
  \item \textbf{Calm}适用人群：压力大的上班族、入睡困难者。核心卖点："每日平静"冥想引导 + 睡眠故事（由名人朗读），界面唯美，英文内容丰富。
  \item \textbf{潮汐}适用人群：需要白噪音专注/入睡的中文用户。核心卖点：国产精品，自然音效 + 冥想课程 + 番茄钟，中文引导更亲切。
  \item \textbf{小睡眠}适用人群：需要助眠音效、记录梦话/呼噜声的人。核心卖点：海量白噪音可叠加组合（雨声+风扇+篝火），睡前定时关闭，梦话录音有趣又实用。
\end{itemize}

\booksubsection{女性健康}

\begin{itemize}
  \item \textbf{Flo}适用人群：需要记录经期、备孕或关注女性健康的用户。核心卖点：AI 预测下次经期/排卵日准确度高，社区可匿名提问，界面设计精致。
  \item \textbf{美柚}适用人群：国内女性用户、希望看中文社区与科普的人。核心卖点：本地化内容丰富，按经期/备孕/孕期/育儿分阶段提供对应知识。
\end{itemize}

\booksubsection{用药提醒}

\begin{itemize}
  \item \textbf{用药助手}适用人群：需要长期服药的慢性病患者、需要照顾老人用药的子女。核心卖点：扫码识别药品、自动建立用药清单，按时弹窗提醒，记录服药情况可分享给家属。
\end{itemize}

\booksubsection{数据中枢}

\begin{itemize}
  \item \textbf{苹果健康/Google Fit}适用人群：所有智能手机用户。核心卖点：汇聚来自手机、手环及第三方健康 App 的所有数据（步数、心率、睡眠、经期等），是统一的"健康档案中心"。
\end{itemize}

\booksection{三、旅行出行}

出门在外，几款趁手的 App 能让你少踩无数坑。下面按出行全流程推荐。

\booksubsection{航班高铁}

\begin{itemize}
  \item \textbf{航旅纵横}适用人群：经常坐飞机的商务人士。核心卖点：官方民航数据，航班动态/登机口变更/延误预警最权威，可查机型与座位图。
  \item \textbf{12306}适用人群：所有国内火车出行者。核心卖点：官方唯一可信购票渠道，支持候补购票（有人退票自动抢），学生票/团体票办理方便。
  \item \textbf{飞常准}适用人群：对航班准点率敏感、需要接送机的人。核心卖点：航班延误预测准确，可分享行程给家人实时查看你的位置与到达时间。
\end{itemize}

\booksubsection{地图导航}

\begin{itemize}
  \item \textbf{高德地图}适用人群：国内城市用户、自驾/打车族。核心卖点：实时路况准确、避开拥堵能力强，打车/打车合乘一站搞定，公共交通换乘方案丰富。
  \item \textbf{百度地图}适用人群：对地点信息、室内导航有需求的用户。核心卖点：全景地图/AR 实景导航，景区/商场室内地图细致，语音包丰富。
  \item \textbf{Google Maps}适用人群：出国旅行者。核心卖点：全球覆盖最全、本地评价真实，街景与离线地图是异国他乡的救星。
\end{itemize}

\booksubsection{翻译}

\begin{itemize}
  \item \textbf{谷歌翻译}适用人群：出国旅行者、跨语言学习者。核心卖点："对话"实时双向翻译、"相机"即时取词翻译是旅行神器，支持离线语言包下载。
\end{itemize}

\booksubsection{汇率换算}

\begin{itemize}
  \item \textbf{极简汇率}适用人群：出国旅行、海淘、跨境购物的人。核心卖点：界面极致清爽，支持 160+ 货币，汇率实时更新，可同时对比多种货币，长按数字快速换算。
\end{itemize}

\booksubsection{天气}

\begin{itemize}
  \item \textbf{系统天气 / 中国天气网}适用人群：所有人，查预报与官方气象预警的首选。核心卖点：厂商天气 App 与「中国天气」等官方渠道对接气象台预警，暴雨、台风、高温等预警会以横幅/通知推送，权威性高。
  \item \textbf{彩云天气}适用人群：需要精准到分钟级降雨预报的人。核心卖点："未来两小时是否下雨"准确度业内领先，降雨带动态地图直观可见，下雨前主动推送提醒。
  \item \textbf{Carrot Weather}适用人群：追求趣味性与丰富 widget 的天气发烧友。核心卖点：毒舌 AI 语音播报天气、超多自定义 widget 桌面组件、雷达图专业，付费版功能更全。
  \item \textbf{预警要点}：天气类 App 要\textbf{允许通知}才收得到预警；预警分蓝、黄、橙、红四级，红色表示灾害性天气风险最高，应按提示调整出行。更多开启方式见本章第七节「天气预警 / 自然灾害预警 / 地震预警」。
\end{itemize}

\booksubsection{旅行攻略}

\begin{itemize}
  \item \textbf{马蜂窝}适用人群：喜欢做详细攻略、自驾游/深度游的人。核心卖点：用户原创游记内容丰富，目的地玩法/路线/避坑指南详尽，可下载离线攻略。
  \item \textbf{小红书}适用人群：想看真实体验、找小众打卡点/美食的年轻旅行者。核心卖点：图文短视频形式直观，"XX 城市 3 天攻略"等内容更新快、时效性强，搜索体验好。
\end{itemize}

\booksection{四、自动化工具启蒙}

自动化是“用一次配置换长期省事”的杠杆，能把你每天重复的几十次操作压缩成零。本节只做概念启蒙，具体玩法与详细教程请见第二十三章：编程思维与自动化启蒙。

\begin{itemize}
  \item \textbf{iOS"快捷指令"}：苹果官方自动化工具，无需编程，可一键完成"到家开 WiFi+亮度调低+播报天气"等组合动作。
  \item \textbf{安卓自动化}：系统级场景联动（如小米"自动任务"、华为"情景智能"）上手最简单；Tasker 功能强大但学习曲线陡峭，适合极客。
  \item \textbf{核心思路}：先识别"每天都要重复做的操作"，再想"能否用一个触发条件 + 多个动作替代"。详细示例与进阶玩法见第 23 章。
\end{itemize}

\booksection{五、浏览器与搜索引擎：你的信息入口}

浏览器是你通往互联网的窗口，搜索引擎是你在信息海洋中导航的罗盘。用好它们，能让你少走很多弯路。

\booksubsection{浏览器的工作原理：}

当你在浏览器中输入网址后，浏览器会：(1)向DNS服务器查询域名对应的IP地址；(2)向该IP地址的服务器发送HTTP/HTTPS请求；(3)服务器返回网页代码（HTML、CSS、JavaScript）；(4)浏览器解析代码并渲染成你看到的页面。

\booksubsection{主流浏览器推荐：}

\begin{itemize}
  \item \textbf{微软Edge}：基于Chromium内核，速度快、兼容好。内置Copilot AI助手，可以总结网页、翻译、对话。支持跨设备同步，个人强烈推荐作为主力浏览器。
  \item \textbf{Safari}：苹果设备默认浏览器，与系统深度集成，省电高效。隐私保护好，但扩展插件较少。
  \item \textbf{Chrome}：全球使用率最高，扩展丰富。但内存占用大，手机端相对耗电。
\end{itemize}

\booksubsection{搜索引擎选择：}

\begin{itemize}
  \item \textbf{Bing}：微软旗下搜索引擎，与Edge深度集成。AI摘要功能可以直接给出答案而不只是链接，搜索效率很高。图片搜索体验也不错。
  \item \textbf{Yandex}：俄罗斯最大的搜索引擎，图片搜索和反向图片搜索能力很强。找相似图片、识别物体时非常好用。
  \item \textbf{百度}：国内信息搜索的主力，但广告较多，需要一定的辨别能力。
  \item \textbf{Google}：全球最精准的搜索引擎，但国内需要特殊网络才能使用。
\end{itemize}

\booksubsection{高效信息检索指南：}

\begin{itemize}
  \item \textbf{善用AI搜索}：Bing的AI摘要、Edge的Copilot可以直接回答问题，帮你从海量结果中提炼关键信息，比逐条点击链接高效得多。
  \item \textbf{识别虚假信息与广告}：搜索结果中标注"广告"或"赞助"的条目是付费推广，不是自然搜索结果。注意甄别信息来源的权威性，优先选择官方网站、知名媒体的内容。
  \item \textbf{正确下载应用}：在搜索引擎中输入"应用名+官网"，如"微信官网"，优先点击带有官方域名标识的结果。不要从第三方下载站安装软件，它们经常捆绑流氓软件。
  \item \textbf{官网鉴别技巧}：下载Windows 11应该访问microsoft.com，下载Chrome应该访问google.com。看域名是否与品牌一致，是最基本的鉴别方法。
\end{itemize}

\booksubsection{网站域名知识：}

域名是网站的“门牌号”，读懂它能帮你判断网站的性质和来源：

\begin{itemize}
  \item \textbf{常见后缀含义}：
  \begin{itemize}
    \item .com — 商业机构，最常见
    \item .cn — 中国国家域名
    \item .org — 非营利组织
    \item .net — 网络服务机构
    \item .edu — 教育机构
    \item .gov — 政府机构
    \item .xyz/.tech/.info — 新兴通用域名，门槛低，需谨慎
  \end{itemize}
  \item \textbf{国家/地区域名}：
  \begin{itemize}
    \item .jp — 日本
    \item .kr — 韩国
    \item .uk — 英国
    \item .de — 德国
    \item .tw — 台湾
    \item .hk — 香港
  \end{itemize}
  \item \textbf{安全提示}：.gov和.edu后缀的网站通常更可信，因为申请门槛高。.xyz、.top等廉价域名容易被钓鱼网站利用，访问时要多留个心眼。
\end{itemize}

\booksubsection{广告拦截：}

\begin{itemize}
  \item \textbf{AdGuard}适用人群：被各类弹窗、贴片广告、应用内广告困扰的所有用户。核心卖点：系统级广告拦截（不仅限浏览器），可过滤 App 内横幅广告与开屏追踪，省流量、加速页面加载，付费版支持全设备。
\end{itemize}

\booksubsection{隐私浏览器：}

\begin{itemize}
  \item \textbf{Brave}适用人群：注重隐私、不希望被广告追踪的日常浏览者。核心卖点：基于 Chromium 内核，开箱即用屏蔽广告与追踪器，自带 Tor 隐身模式，速度比 Chrome 快、内存占用更低。
  \item \textbf{DuckDuckGo}适用人群：希望"搜一个词不被全网追着推广告"的用户。核心卖点：自带同名搜索引擎（不记录搜索历史），一键清除所有浏览痕迹（"燃烧按钮"），跨平台体验统一。
\end{itemize}

\booksubsection{浏览器对比一览：}

\begin{center}
\begin{tabularx}{\linewidth}{|X|X|X|}
\hline
\textbf{浏览器} & \textbf{核心卖点} & \textbf{适用人群} \tabularnewline
\hline
微软 Edge & Chromium 内核 + 内置 Copilot AI，可总结网页/翻译/对话 & 追求效率的办公族、AI 重度用户（推荐主力） \tabularnewline
Safari & 与苹果生态深度集成，省电高效，隐私保护强 & iPhone/Mac 用户，重视续航与系统联动 \tabularnewline
Chrome & 全球使用率最高，扩展插件生态最丰富 & 需要大量扩展插件的进阶用户、跨设备同步党 \tabularnewline
Brave & 开箱即用屏蔽广告与追踪器，速度更快 & 注重隐私、不装插件也想屏蔽广告的用户 \tabularnewline
DuckDuckGo & 同名搜索引擎不记录历史，一键清除痕迹 & 不希望被个性化追踪的隐私敏感用户 \tabularnewline
\hline
\end{tabularx}
\end{center}

\begin{tcbpractice}{试着创建一个"睡前放松"的自动化流程}
\begin{enumerate}
  \item 创建触发条件：特定时间（如22:30）。
  \item 添加执行动作：调低音量 > 打开勿扰模式 > 播放放松歌单 > 调低亮度。
  \item 保存并测试。
\end{enumerate}

\end{tcbpractice}

\booksection{六、实用系统工具：清理、防骚扰与安全防护}

除了效率与健康类 App，真正影响日常体验的往往是「少被打扰、存储够用、骗子进不来」这三件事。本节先盘点常用实用软件，再单独展开\textbf{手机自带的骚扰拦截功能}——很多人买了手机多年，却从没打开过这些开关。

\booksubsection{系统清理与存储优化}

\begin{itemize}
  \item \textbf{清浊}适用人群：安卓用户，希望轻量清理缓存、残留文件、空文件夹，又不想被「手机管家」全家桶打扰的人。核心卖点：体积小、无广告干扰，支持深度扫描应用残留、大文件定位、存储占用分析；比商业「一键加速」更克制，适合动手能力强的用户做定期维护。注意：清理前先确认路径，微信、相册、文档目录不要盲目全选；系统级缓存尽量优先用品牌自带存储清理，第三方工具只作补充。
  \item \textbf{系统自带存储清理}适用人群：所有人，尤其是长辈和怕折腾的用户。核心卖点：安卓 [设置] > [存储] > [清理加速]；iOS [设置] > [通用] > [iPhone 储存空间]。系统更懂自家 App 缓存结构，误删风险最低，应作为\textbf{第一选择}。
  \item \textbf{微信/QQ 内置存储管理}适用人群：聊天大户。核心卖点：微信 [我] > [设置] > [通用] > [存储空间]，可按聊天、文件类型精细清理；比外部「清微信」类工具更安全。
  \item \textbf{原则}：能用系统/App 自带的，就不用第三方；能按路径/类型勾选的，就不要「一键全清」。清理的目标是释放空间，不是把重要数据一起扫掉。
\end{itemize}

\booksubsection{防骚扰与安全类实用软件}

\begin{itemize}
  \item \textbf{国家反诈中心}适用人群：所有人，尤其是长辈、财务相关岗位、经常网购/接陌生电话的人。核心卖点：官方出品，可开启来电预警、短信风险提示、可疑网址检测；与公安反诈数据库联动，是「三道防线」里最有官方背书的一环。
  \item \textbf{品牌安全中心 / 手机管家}适用人群：对应品牌用户。核心卖点：华为手机管家、小米安全中心、OPPO 手机管家、vivo 手机管家等，集成骚扰拦截、病毒查杀、权限管理、支付保护，通常\textbf{免费且与系统底层协同更好}，不建议再装第二套商业安全软件造成权限冲突。
  \item \textbf{运营商防骚扰服务}适用人群：被高频外呼、AI 语音电话轰炸的用户。核心卖点：移动发送KTFSR到 10086 开通；电信关注「天翼防骚扰」；联通在「联通手机管家」中设置。拦截发生在\textbf{运营商侧}，手机不响铃也能挡住一批骚扰号。
  \item \textbf{AdGuard}适用人群：被各类弹窗、贴片广告、应用内广告困扰的所有用户。核心卖点：系统级广告拦截（不仅限浏览器），可过滤 App 内横幅广告与开屏追踪，省流量、加速页面加载，付费版支持全设备。
  \item \textbf{慎装「清理加速 / 骚扰拦截」全家桶}：部分第三方管家常驻后台、反复请求通讯录与短信权限，拦截能力未必比系统强，反而增加耗电与隐私风险。优先打开系统自带功能，再考虑官方反诈 App。
\end{itemize}

\textbf{实用软件选型原则}\par

\begin{itemize}
  \item \textbf{系统优先}：拦截、清理、权限管理，先打开手机自带开关。
  \item \textbf{官方补充}：国家反诈中心、运营商防骚扰属于官方通道，可叠加。
  \item \textbf{第三方克制}：同一类功能装一个就够，避免权限打架与后台抢资源。
  \item \textbf{少即是多}：工具装得越多，通知和权限面越大，被「工具本身」打扰的概率越高。
\end{itemize}

\booksubsection{手机自带骚扰拦截功能：各品牌怎么开}

下面按品牌列出入口。名称可能随系统版本略有差异，可在设置里直接搜索「骚扰」「拦截」「未知来电」。与第八章「骚扰电话治理」互为补充：这里侧重\textbf{开关位置}，那里侧重\textbf{号码标记与误标申诉}。

\begin{center}
\begin{tabularx}{\linewidth}{|X|X|X|}
\hline
\textbf{平台} & \textbf{入口路径} & \textbf{推荐开启项} \tabularnewline
\hline
iPhone / iOS & {}[设置] > [电话] / [信息] & 静音未知来电；已阻止的联系人；信息中「过滤未知发件人」 \tabularnewline
小米 / Redmi & 电话 App > 设置 > 骚扰拦截；或 安全中心 & 智能拦截、陌生号响一声挂断、短信关键词过滤 \tabularnewline
华为 / 荣耀 & 电话 > 骚扰拦截；手机管家 > 骚扰拦截 & 来电识别、黑名单、短信智能拦截 \tabularnewline
OPPO / realme / 一加 & 电话 App > 骚扰拦截；或 本机 > 骚扰拦截 & 来电拦截、诈骗电话识别、短信拦截规则 \tabularnewline
vivo / iQOO & 电话 > 骚扰电话拦截；i 管家 & 骚扰拦截、诈骗电话预警、信息过滤 \tabularnewline
三星 & 电话 > 设置 > 阻止号码 / 来电显示与垃圾电话 & 阻止未知号码、垃圾电话识别（地区支持度不同） \tabularnewline
\hline
\end{tabularx}
\end{center}

\begin{tcbinfo}{三层防护叠加建议}
日常推荐「\textbf{系统拦截 + 运营商防骚扰 + 国家反诈中心}」三层叠加：系统管标记与黑名单，运营商在网侧挡高频外呼，反诈中心补官方预警。三层不必都开到最严，避免把快递、外卖、医生电话也误杀——可对通讯录与最近联系人放行。

\end{tcbinfo}

\begin{tcbpractice}{动手练习：今天就把骚扰拦截打开}
\begin{enumerate}
  \item 打开电话 App 设置，找到「骚扰拦截 / 静音未知来电」，开启智能拦截。
  \item 给爸妈手机也检查一遍：同样路径 + 开通运营商防骚扰（移动可发 KTFSR 到 10086）。
  \item 应用商店安装「国家反诈中心」，注册并打开来电预警。
  \item 试打一个已被标记的测试号（或看拦截记录），确认功能生效。
\end{enumerate}

\end{tcbpractice}

\booksection{七、天气预警、自然灾害预警与地震预警}

骚扰拦截挡住的是「电话诈骗」，预警功能挡住的是「天灾与极端天气」。很多人只装了天气 App 看温度，却\textbf{没开通知、没开地震预警}——暴雨红色预警、台风、地震来临前那几秒到几十秒，手机本可以帮你争取到宝贵的避险时间。

\booksubsection{天气预警：暴雨、台风、高温怎么收}

\begin{itemize}
  \item \textbf{官方渠道优先}：中国气象局 / 中央气象台 / 各省气象局发布的预警最权威。可关注「中国天气」微信公众号、微博，或下载「中国天气」等官方类 App。
  \item \textbf{系统天气 App}：华为天气、小米天气、OPPO 天气、vivo 天气、苹果「天气」等，多数会推送当地气象台预警；进入天气 App 设置，确认\textbf{「气象预警 / 灾害性天气」通知已开启}。
  \item \textbf{彩云天气等分钟级降雨}：适合通勤、接送孩子前判断「这半小时要不要带伞」；预警能力仍以官方气象台为准。
  \item \textbf{预警颜色含义（通用）}：
  \begin{itemize}
    \item \textbf{蓝色}：注意防范，一般照常出行、关注更新。
    \item \textbf{黄色}：加强防范，减少不必要户外活动。
    \item \textbf{橙色}：严重，尽量留在室内，远离危险区域。
    \item \textbf{红色}：特别严重，听从停课、停工、避险指引，不要硬闯积水、山洪、大风路段。
  \end{itemize}
  \item \textbf{常见预警类型}：暴雨、暴雪、台风、高温、寒潮、大风、雷电、冰雹、道路结冰、大雾、干旱、地质灾害风险。
\end{itemize}

\booksubsection{自然灾害预警：不只「下雨」}

\begin{itemize}
  \item \textbf{发布主体}：气象台负责气象类；地质灾害、山洪等常由自然资源、水利、应急管理部门联合发布；地方政府也会通过应急广播、短信、12379 等渠道发布。
  \item \textbf{怎么多渠道接收}：
  \begin{itemize}
    \item 手机系统通知 + 天气 App 预警。
    \item 微信关注当地「应急管理局」「气象局」公众号。
    \item 12379 国家预警信息发布相关渠道（短信/服务，视地区开通情况）。
    \item 电视、应急广播、村/社区大喇叭——山区、汛期不要只盯手机。
  \end{itemize}
  \item \textbf{户外与出行提醒}：台风天不贴广告牌、不走树下；暴雨不进地下车库、隧道、下穿桥；山区旅游遇「地质灾害气象风险」预警应立即离开陡坡、沟谷。
  \item \textbf{与急救手册联动}：灾后失联、手机没电、需要定位求救时，见第三十二章地震/灾害场景，以及第七章卫星通信与无网通。
\end{itemize}

\booksubsection{地震预警：和「地震预报」不是一回事}

\textbf{地震预警}（EEW）是在地震波（尤其是破坏力更强的横波）到达目标区之前，利用电磁波比地震波快、以及台站更快捕捉到纵波的原理，向可能受影响区域\textbf{提前数秒到数十秒}推送警报。它\textbf{不能预测地震是否发生}，但能抢出「就近避险、切断火源、冲出电梯」的窗口。距离震中越近，预警时间越短，甚至可能来不及。

\begin{itemize}
  \item \textbf{官方与行业方案}：中国大陆常见接入方包括中国地震预警网相关能力，以及成都高新减灾研究所等机构推动的预警服务；部分厂商手机会预装或可联网接收预警。
  \item \textbf{第三方 App}：应用商店搜索「地震预警」等官方/机构相关应用（安装前看清开发者，避免山寨「假预警」App 要你开一堆无关权限）。
  \item \textbf{收到警报后怎么做}：室内就近躲避（桌下、承重墙角、坚固家具旁），护住头颈；远离窗户、外墙、高大家具；不要盲目冲楼梯/跳楼；电梯里立刻在最近楼层离开；室外到开阔地，远离电线、广告牌、危房。震后检查燃气与电路，注意余震。
  \item \textbf{勿扰勿关}：地震预警应允许「突破静音/勿扰」响铃；睡前把手机放伸手可及处，铃声音量别调到 0。
\end{itemize}

\booksubsection{手机自带地震预警：各品牌怎么开}

下列入口随系统版本可能略有差异，可在设置中直接搜索「地震」「预警」「安全预警」。iOS 系统本身对国内地震预警支持有限，通常需安装可靠的第三方预警 App，并保持定位与通知权限。

\begin{center}
\begin{tabularx}{\linewidth}{|X|X|X|}
\hline
\textbf{平台} & \textbf{常见入口} & \textbf{说明} \tabularnewline
\hline
小米 / Redmi & 手机管家 / 安全中心 > 地震预警；或 设置搜索「地震」 & 部分版本系统内置预警服务，需联网并开启定位 \tabularnewline
华为 / 荣耀 & 设置 > 安全 > 应急预警通知 / 地震预警 & 名称可能是「应急预警通知」，内含地震等 \tabularnewline
OPPO / realme / 一加 & 设置 > 安全（SOS 紧急联络）> 地震预警 & 部分机型与地区服务覆盖有关 \tabularnewline
vivo / iQOO & 天气 App 或 设置 > 安全/更多 > 地震预警 & 以实际系统版本菜单为准 \tabularnewline
iPhone / iOS & 系统设置中通常无完整内置开关；安装可靠地震预警 App 并允许通知 & App 需后台运行与定位权限；系统勿扰勿屏蔽其警报 \tabularnewline
三星等 & 设置 > 安全与紧急情况 / 无线紧急警报 & 「紧急警报」覆盖范围因地区与运营商而异 \tabularnewline
\hline
\end{tabularx}
\end{center}

\begin{tcbwarn}{关于地震预警的常见误区}
\begin{itemize}
  \item \textbf{预警 $\neq$ 预报}：它不能告诉你「明天会不会地震」，只在已发生地震、波尚未到达时报警。
  \item \textbf{不是哪里都能提前很多秒}：震中附近可能只有数秒甚至来不及，预警主要惠及周边区域。
  \item \textbf{关通知就等于没装}：再好的预警，静音、杀后台、不给定位，都可能收不到。
  \item \textbf{山寨 App}：警惕要通讯录、短信权限的「地震预报」软件，优先系统内置或正规机构方案。
\end{itemize}

\end{tcbwarn}

\begin{tcbinfo}{预警能力怎么搭配}
日常建议：「\textbf{系统天气预警通知 + 官方气象/应急渠道 + 地震预警（系统或正规 App）}」。与防骚扰一样，不必装一堆同类软件；关键是\textbf{通知权限打开、定位允许、勿扰能突破}。给长辈手机做完骚扰拦截后，顺手把天气预警和地震预警也检查一遍。

\end{tcbinfo}

\begin{tcbpractice}{动手练习：10 分钟检查预警设置}
\begin{enumerate}
  \item 打开系统天气 App，进入设置，开启「气象预警 / 灾害性天气」通知。
  \item 在系统设置里搜索「地震」或「预警」，找到并开启地震预警/应急预警（安卓）；iPhone 则安装正规地震预警 App 并允许通知。
  \item 关注当地气象局或应急局的公众号/服务号，允许其推送。
  \item 确认预警声音不会被勿扰完全静音；给爸妈手机同样检查一遍。
\end{enumerate}

\end{tcbpractice}

\booksection{八、跨设备协同：远控电脑、快传与账号安全}

手机用得顺，多半是因为它和电脑、平板「组了队」。本节补充几类高频实用工具：手机远控电脑、AI 办公与远程协作、动态壁纸、局域网快传、跨设备剪贴板、动态口令 2FA，以及纸质扫描与日程任务的衔接。更系统的多设备联动见第十七章。

\booksubsection{手机远控电脑：向日葵 / ToDesk}

\begin{itemize}
  \item \textbf{向日葵远程控制}适用人群：需要在外出时操作家里/公司 Windows 或 Mac 的用户；也常用于帮长辈远程修手机（见第二十七章）。核心卖点：国产生态成熟，手机 App 控制电脑桌面、传文件、开关机（部分需硬件）；免费版一般够个人轻度使用。
  \item \textbf{ToDesk}适用人群：追求连接稳定、低延迟远程办公/运维的用户。核心卖点：设备码+密码连接，手机控电脑、电脑控手机（视版本）、文件传输与多显示器切换较顺手，国内节点体验通常较好。
  \item \textbf{基本用法}：电脑端安装客户端并登录/保持在线 手机端装 App 用\textbf{设备识别码 + 连接密码}（或同一账号）建立会话 手机触控板/手势操作桌面。
  \item \textbf{安全红线（务必记住）}：
  \begin{itemize}
    \item 陌生人自称「客服/公检法/运营商」让你装向日葵、ToDesk 并远程你\textbf{手机或电脑}，\textbf{一律拒绝}——这是典型诈骗手法（详见第二十七章）。
    \item 自己用时：开启二次验证、修改默认密码、不用时退出账号或关闭无人值守；公司电脑遵守单位安全规定。
    \item 公共场合远程办公注意周围是否有人偷窥屏幕与口令。
  \end{itemize}
\end{itemize}

\booksubsection{AI 办公软件：Trae / WorkBuddy / Dumate}

\begin{itemize}
  \item \textbf{Trae}适用人群：编程与文档协作用户，中文环境友好。核心卖点：字节跳动推出的 AI IDE/编程助手，SOLO 模式可一句话生成项目；日常写代码、改文档效率高。若电脑端运行，可在移动端配合远控或官方配套能力查看/跟进任务（以产品最新形态为准）。详见第四十一章AI 编程工具全景。
  \item \textbf{WorkBuddy}适用人群：希望 AI 帮忙整理会议纪要、邮件、汇报、办公文档的职场用户。核心卖点：面向办公场景的 AI 助手，适合「起草—润色—总结」这类高频文案；部分工作流可在手机端发起，重活放到电脑端完成。
  \item \textbf{Dumate}适用人群：需要多人协作、文档与知识管理一体化的团队/个人。核心卖点：AI 协作向办公工具，强调资料汇聚与协同编辑；手机端可查看与轻量编辑，复杂排版建议在电脑大屏处理。
  \item \textbf{「手机端远程控制电脑」怎么理解}：很多 AI 办公能力跑在电脑/网页端；在手机上通过\textbf{向日葵/ToDesk}连回电脑，或使用产品自身的手机 App/网页，就能在通勤时「遥控」电脑上的 AI 任务。原则：\textbf{重计算、长文档放电脑；看进度、改几句放手机}。
  \item \textbf{隐私提醒}：上传合同、身份证、客户数据前看清产品隐私条款与企业合规要求；公司资料优先用单位指定工具。
\end{itemize}

\booksubsection{人工桌面：米哈游动态壁纸（手机 + PC）}

\begin{itemize}
  \item \textbf{人工桌面}适用人群：喜欢米哈游角色、想要免费动态壁纸的用户。核心卖点：米哈游推出的\textbf{免费}动态壁纸应用，支持\textbf{手机端与 PC 端}；角色 Live2D/动态演出可作桌面，部分机型/系统可设为锁屏或桌面壁纸。
  \item \textbf{省电与性能}：动态壁纸比静态图更耗电、占 GPU；低电量或老机型可只在充电时使用，或降低帧率（若 App 提供选项）。
  \item \textbf{安装来源}：优先官网/正规应用商店，避免第三方「破解版」捆绑木马；PC 版注意安装路径与开机自启是否需要。
\end{itemize}

\booksubsection{局域网文件快传：LocalSend / 飞鸽传书 / 隔空投送}

\begin{itemize}
  \item \textbf{LocalSend}适用人群：安卓/Windows/iOS/Mac/Linux 混用，想要\textbf{跨平台、不传公有云}快传的人。核心卖点：开源、局域网直连，同一 Wi-Fi 下互传照片视频文档，无需数据线、不经过厂商服务器，速度取决于路由器。
  \item \textbf{飞鸽传书}适用人群：办公室 Windows 电脑之间、老局域网办公场景。核心卖点：老牌局域网传输/聊天工具，适合内网传大文件；注意从正规渠道下载，企业内网遵守 IT 策略。
  \item \textbf{隔空投送（AirDrop）}适用人群：iPhone、iPad、Mac 用户。核心卖点：苹果生态一键分享，蓝牙发现 + Wi-Fi 直连传输，免流量；可设「仅联系人」降低陌生人骚扰。
  \item \textbf{其他补充}：小米互传、华为分享、vivo 互传等品牌方案在同一品牌设备间很快；微信「文件传输助手」方便但受文件大小与登录设备限制，大文件更推荐 LocalSend/隔空投送。
  \item \textbf{使用提示}：双方连\textbf{同一 Wi-Fi}；公共 Wi-Fi 传敏感文件有风险，优先私人热点；传输完成后在接收方确认文件是否完整。
\end{itemize}

\booksubsection{跨设备剪贴板同步：KDE Connect / 手机电脑互传}

\begin{itemize}
  \item \textbf{苹果「通用剪贴板」}适用人群：iPhone + Mac 用户。核心卖点：登录同一 Apple ID、开接力后，手机复制 Mac 直接粘贴（反之亦然），文字与部分文件均可。
  \item \textbf{KDE Connect}适用人群：安卓 + Windows/Linux/Mac，想要开源剪贴板同步、通知镜像、远程输入的用户。核心卖点：局域网配对后同步剪贴板、推送手机通知到电脑、用手机当幻灯片翻页器/触控板；数据不经过第三方账号。
  \item \textbf{品牌「手机电脑互传」}：华为多屏协同/超级终端、小米妙享、vivo 互传、OPPO 互连、荣耀互联等，常自带\textbf{跨屏剪贴板与文件拖拽}，同品牌体验最好。
  \item \textbf{Windows + 安卓官方通道}：「手机连接 / Link to Windows」（Your Phone 演进）可在电脑上看手机通知、照片，并在支持的场景复制粘贴；三星等机型集成更深。
  \item \textbf{安全提醒}：剪贴板可能残留密码、验证码；在共用电脑上慎用全局同步，用完在密码框避免二次粘贴泄露；部分输入法也会读剪贴板，注意权限。
\end{itemize}

\booksubsection{动态口令 2FA：Aegis / Microsoft Authenticator}

\begin{itemize}
  \item \textbf{为什么要用认证器而不是短信}：短信验证码可能遭遇 SIM 补卡劫持；TOTP 动态口令在本地按时间算法生成，无需依赖手机号信号，更抗钓鱼与拦截（仍需防木马与截屏）。
  \item \textbf{Aegis Authenticator}适用人群：安卓用户，偏好开源、本地可控。核心卖点：开源免费，支持加密备份与导入导出，界面干净，可指纹/人脸解锁令牌列表。
  \item \textbf{Microsoft Authenticator}适用人群：使用微软账号、Office 365/企业账号的用户；跨 iOS/安卓。核心卖点：推送批准登录、密码监控与账号备份（按产品版本），与 Windows/企业身份体系结合紧密。
  \item \textbf{Google Authenticator / 2FAS 等}：生态内也常见；关键是有\textbf{加密备份}，换机不丢令牌。
  \item \textbf{换机必做}：旧机迁移或用服务里的「备用码/恢复密钥」；\textbf{备用码打印在纸上}存两处，不要只存同一部手机。
  \item \textbf{与密码管理器}：部分密码库也可存 TOTP，方便但风险更集中；高价值账号建议「密码管理器 + 独立认证器」分工。详见第十四章。
\end{itemize}

\booksubsection{推荐组合：一天里怎么串起来}

\textbf{实用工作流示例}\par

\begin{itemize}
  \item \textbf{早上}：滴答清单看今日日程 天气/预警扫一眼 电脑开 Trae/文档，手机放一旁。
  \item \textbf{白天}：纸质材料用扫描类 App 转 PDF，经 LocalSend/隔空投送发到电脑；长文在电脑用 AI 办公工具起草。
  \item \textbf{外出}：用 ToDesk/向日葵连回电脑继续任务；重要账号登录走 Aegis/Microsoft Authenticator 动态口令。
  \item \textbf{晚上}：滴答清单勾选完成项；需要时人工桌面换个壁纸放松一下。
\end{itemize}

\begin{tcbwarn}{远控与账号安全口诀}
\begin{itemize}
  \item 别人让你装远控软件\textbf{挂电话，自己找官方客服}。
  \item 验证码、动态口令\textbf{不告诉任何人}，包括「快递」「客服」。
  \item 公司电脑远控与安装 AI 软件前，先问 IT 能不能用。
  \item 公共网络传文件、远程登录时格外警惕。
\end{itemize}

\end{tcbwarn}

\begin{tcbpractice}{动手练习：把「跨设备」链路跑通一次}
\begin{enumerate}
  \item 手机与电脑连同一 Wi-Fi，用 LocalSend（或隔空投送）互传一张照片。
  \item 在滴答清单建一条「今晚 8 点：整理扫描件」任务，确认手机与电脑都能看到。
  \item 为邮箱或微信开启 2FA（若支持），用 Aegis 或 Microsoft Authenticator 绑定，并保存备用码。
  \item 在自己电脑上安装向日葵或 ToDesk，用自己手机成功连一次，再\textbf{退出或加固密码}，熟悉流程但不留安全后门。
\end{enumerate}

\end{tcbpractice}
